%% kapa: exam=1 round=27 year=2016 period=1 form=A subject=civil_law subject_ko=민법 seq=33 number=33
%% source: Q-Net 감정평가사 기출문제 (cst00309) artlSeq=5169402 file=extracted/1008/5169402/제27회 감정평가사 1차 문제/제27회 감정평가사 1교시문제(A형).pdf
%% conversion: kordoc 4.24.51 -> tools/gen1.py (auto); review=proofread-sonnet; flags=none
\begin{question}{33}
건물의 구분소유 및 집합건물 등에 관한 설명으로 옳지 않은 것은? (다툼이 있으면 판례에 따름)
\begin{choices}
  \item 공용부분을 전유부분으로 변경하기 위하여는 구조상으로나 이용상으로 다른 전유부분과 독립되어 있어야 한다.
  \item 구분소유자 중 일부가 복도, 계단과 같은 공용부분의 일부를 아무런 권원 없이 점유ㆍ사용하는 경우, 특별한 사정이 없는 한 다른 구분소유자들에게 임료 상당의 손해가 발생한 것으로 볼 수 있다.
  \item 대지에 대한 지상권도 대지사용권이 될 수 있다.
  \item 집합건물의 관리단은 구분소유자 전원을 구성원으로 하며, 별도의 설립행위가 필요하지 않다.
  \item 구분건물이 물리적으로 완성되기 전이라도 건축허가신청 등을 통하여 구분의사가 객관적으로 표시되면 구분행위의 존재를 인정할 수 있다.
\end{choices}
\end{question}
\answer{2}
